\section{Цель работы}

Ознакомление с основными методами обработки событий и прерываний в микроконтроллерах.

\section{Постановка задачи}

Разработать программу для микроконтроллера STM32F200, обеспечивающую мигание светодиода PG7, а также обработку нажатий кнопок WAKEUP и USER с использованием прерываний.
При нажатии кнопки WAKEUP должен включаться светодиод PG6, а при нажатии кнопки USER --- светодиод PG8.
Обработка прерываний должна выполняться с разными приоритетами.

\section{Выполнение работы}

В среде Keil $\mu$Vision5 был создан проект, в котором реализована программа обработки внешних прерываний и управления светодиодами.

\subsection{Функция задержки}

Для создания временных задержек используется функция \verb|delay()|. Она реализована в виде пустого цикла; переменная цикла объявлена с квалификатором \verb|volatile|, чтобы компилятор не оптимизировал цикл.

\begin{listing}[H]
\caption{Функция задержки}
\begin{minted}[linenos]{c}
void delay(unsigned long time)
{
    volatile unsigned long i;
    for (i = 0; i < time; i++) {}
}
\end{minted}
\end{listing}

\subsection{Обработчик прерывания WAKEUP}

Обработчик \verb|EXTI0_IRQHandler| вызывается при нажатии кнопки WAKEUP, подключенной к выводу PA0.
В начале работы обработчика проверяется флаг прерывания в регистре \verb|EXTI->PR|.
После подтверждения события включается светодиод PG6, выполняется задержка, светодиод выключается, а флаг прерывания сбрасывается.

\begin{listing}[H]
\caption{Обработчик прерывания WAKEUP}
\begin{minted}[linenos]{c}
void EXTI0_IRQHandler(void)
{
    if (EXTI->PR & EXTI_PR_PR0) {
        GPIOG->ODR |=  (1 << 6);   /* Включить светодиод PG6 */
        delay(3000000);
        GPIOG->ODR &= ~(1 << 6);   /* Выключить светодиод PG6 */
        EXTI->PR = EXTI_PR_PR0;    /* Сброс флага прерывания */
    }
}
\end{minted}
\end{listing}

\subsection{Обработчик прерывания USER}

Обработчик \verb|EXTI15_10_IRQHandler| срабатывает при нажатии кнопки USER, подключенной к выводу PG15.
Так как данный обработчик обслуживает несколько линий EXTI, дополнительно проверяется, что событие произошло именно на линии 15.
Дальнейшая логика аналогична обработчику WAKEUP: включение светодиода PG8, задержка, выключение и сброс флага прерывания.

\begin{listing}[H]
\caption{Обработчик прерывания USER}
\begin{minted}[linenos]{c}
void EXTI15_10_IRQHandler(void)
{
    if (EXTI->PR & EXTI_PR_PR15) {
        GPIOG->ODR |=  (1 << 8);   /* Включить светодиод PG8 */
        delay(3000000);
        GPIOG->ODR &= ~(1 << 8);   /* Выключить светодиод PG8 */
        EXTI->PR = EXTI_PR_PR15;   /* Сброс флага прерывания */
    }
}
\end{minted}
\end{listing}

\subsection{Настройка тактирования и GPIO}

На первом этапе работы программы включается тактирование портов GPIOA и GPIOG.
Светодиоды PG6, PG7 и PG8 настраиваются как выходы, а кнопки PA0 и PG15 --- как входы.

\begin{listing}[H]
\caption{Настройка тактирования и GPIO}
\begin{minted}[linenos]{c}
/* Включение тактирования портов GPIOA и GPIOG */
RCC->AHB1ENR |= RCC_AHB1ENR_GPIOAEN | RCC_AHB1ENR_GPIOGEN;

/* PG6, PG7, PG8 -> выход */
GPIOG->MODER &= ~(3 << (6 * 2));
GPIOG->MODER |=  (1 << (6 * 2));

GPIOG->MODER &= ~(3 << (7 * 2));
GPIOG->MODER |=  (1 << (7 * 2));

GPIOG->MODER &= ~(3 << (8 * 2));
GPIOG->MODER |=  (1 << (8 * 2));

/* PA0 -> вход (кнопка WAKEUP) */
GPIOA->MODER &= ~(3 << (0 * 2));

/* PG15 -> вход (кнопка USER) */
GPIOG->MODER &= ~(3 << (15 * 2));
\end{minted}
\end{listing}

\subsection{Настройка системы прерываний EXTI}

Для работы внешних прерываний включается модуль SYSCFG и настраиваются линии EXTI.
Линия EXTI0 связывается с выводом PA0, а линия EXTI15 --- с выводом PG15.
После этого разрешается генерация прерываний и задается срабатывание по переднему и заднему фронтам сигнала через регистры \verb|RTSR| и \verb|FTSR|.

\begin{listing}[H]
\caption{Настройка EXTI}
\begin{minted}[linenos]{c}
/* Включить тактирование модуля SYSCFG */
RCC->APB2ENR |= RCC_APB2ENR_SYSCFGEN;

/* Привязать EXTI0 к PA0 и EXTI15 к PG15 */
SYSCFG->EXTICR[0] |= SYSCFG_EXTICR1_EXTI0_PA;
SYSCFG->EXTICR[3] |= SYSCFG_EXTICR4_EXTI15_PG;

/* Разрешить прерывания по линиям 0 и 15 */
EXTI->IMR |= (1 << 0) | (1 << 15);

/* Срабатывание по переднему фронту */
EXTI->RTSR |= (1 << 0) | (1 << 15);

/* Срабатывание по заднему фронту */
EXTI->FTSR |= (1 << 0) | (1 << 15);
\end{minted}
\end{listing}

\subsection{Настройка NVIC и приоритетов}

Контроллер NVIC используется для управления прерываниями и задания их приоритетов.
Прерыванию от кнопки USER назначается более высокий приоритет, чем прерыванию от кнопки WAKEUP.
Это означает, что при одновременном возникновении событий сначала будет обработано прерывание USER.

\begin{listing}[H]
\caption{Настройка NVIC}
\begin{minted}[linenos]{c}
NVIC_SetPriority(EXTI15_10_IRQn, 0); /* USER: высокий приоритет */
NVIC_SetPriority(EXTI0_IRQn,     1); /* WAKEUP: низкий приоритет */

NVIC_EnableIRQ(EXTI0_IRQn);
NVIC_EnableIRQ(EXTI15_10_IRQn);
\end{minted}
\end{listing}

\subsection{Основной цикл программы}

В основном цикле программы реализовано мигание светодиода PG7.
Цикл выполняется бесконечно, однако может быть прерван в любой момент при возникновении внешнего прерывания от одной из кнопок.

\begin{listing}[H]
\caption{Основной цикл программы}
\begin{minted}[linenos]{c}
while (1) {
    GPIOG->ODR |=  (1 << 7);   /* Включить PG7 */
    delay(2000000);
    GPIOG->ODR &= ~(1 << 7);   /* Выключить PG7 */
    delay(2000000);
}
\end{minted}
\end{listing}

\section{Вывод}

В ходе выполнения лабораторной работы была разработана программа для микроконтроллера STM32F200, реализующая обработку внешних прерываний.
Были изучены принципы работы системы EXTI и контроллера NVIC, а также механизм приоритетов прерываний.
Практически продемонстрировано, что прерывания с более высоким приоритетом обрабатываются раньше.
Полученные навыки позволяют использовать прерывания для эффективной обработки внешних событий в микроконтроллерных системах.
